\chapter{Tests et Validation}

\section{Introduction}

Ce chapitre décrit les tests effectués pour valider le bon fonctionnement de LudoHeritage.
Les tests couvrent les aspects fonctionnels (scénarios utilisateur), les aspects techniques
(API backend) et la qualité de l'expérience utilisateur.

% ─────────────────────────────────────────────────────────────────────────────
\section{Tests fonctionnels}
% ─────────────────────────────────────────────────────────────────────────────

\subsection{Scénarios testés}

Les scénarios suivants ont été validés manuellement~:

\begin{table}[H]
\centering
\renewcommand{\arraystretch}{1.4}
\begin{tabularx}{\textwidth}{>{\bfseries}p{1cm} p{5cm} X >{\bfseries}p{2cm}}
\toprule
\textbf{N°} & \textbf{Scénario} & \textbf{Résultat attendu} & \textbf{Statut} \\
\midrule
1  & Inscription avec un e-mail valide
   & Compte créé, redirection vers l'onboarding & Validé \\
2  & Inscription avec un e-mail déjà utilisé
   & Message d'erreur \og{}Ce compte existe déjà\fg{} & Validé \\
3  & Connexion avec mauvais mot de passe
   & Message \og{}Identifiants invalides\fg{} & Validé \\
4  & Mot de passe et confirmation différents
   & Blocage côté frontend, message d'erreur & Validé \\
5  & Filtrage du catalogue par région \og{}Afrique\fg{}
   & Seuls les jeux africains s'affichent & Validé \\
6  & Clic sur un jeu depuis la carte SVG
   & La fiche du jeu s'ouvre immédiatement & Validé \\
7  & Comparaison de deux jeux avec ludèmes communs
   & Les ludèmes partagés apparaissent en vert & Validé \\
8  & Soumission d'un score de quiz
   & Score enregistré en MongoDB, classement mis à jour & Validé \\
9  & Ajout d'un jeu en favori, consultation du profil
   & Le jeu apparaît dans \og{}Mes Favoris\fg{} & Validé \\
10 & Modification des préférences depuis le profil
   & Les nouvelles préférences sont sauvegardées & Validé \\
11 & Consultation de l'Awale jouable
   & Les deux joueurs peuvent semer et capturer & Validé \\
12 & LudoBot avec Ollama hors ligne
   & Réponse de secours cohérente affichée & Validé \\
\bottomrule
\end{tabularx}
\caption{Tableau de validation des scénarios fonctionnels}
\label{tab:tests}
\end{table}

% ─────────────────────────────────────────────────────────────────────────────
\section{Tests de l'API backend}
% ─────────────────────────────────────────────────────────────────────────────

Les \textit{endpoints} de l'API REST ont été testés avec des requêtes HTTP directes~:

\begin{itemize}
    \item \texttt{POST /api/auth/register} -- Retourne 200 avec l'utilisateur créé,
          ou 409 (Conflict) si l'e-mail existe déjà.
    \item \texttt{GET /api/jeux} -- Retourne les 20 jeux, même si MongoDB est hors ligne
          (grâce aux données de secours en mémoire).
    \item \texttt{GET /api/jeux/day} -- Retourne un jeu différent chaque jour selon le
          numéro du jour de l'année.
    \item \texttt{POST /api/quiz/scores} -- Enregistre le score et retourne l'objet créé.
    \item \texttt{GET /api/quiz/leaderboard} -- Retourne les dix meilleurs scores triés.
    \item \texttt{GET /api/agent/ping} -- Retourne \texttt{\{"status":"ok"\}} si Ollama répond,
          ou \texttt{\{"status":"error"\}} sinon.
\end{itemize}

% ─────────────────────────────────────────────────────────────────────────────
\section{Gestion de la sécurité}
% ─────────────────────────────────────────────────────────────────────────────

Plusieurs mécanismes de sécurité ont été validés~:

\begin{enumerate}
    \item \textbf{Hashage des mots de passe}~: aucun mot de passe n'est stocké en clair.
          BCrypt génère un hash différent pour chaque utilisateur.

    \item \textbf{Validation des entrées}~: les annotations Jakarta Validation (\texttt{@NotBlank},
          \texttt{@Email}, \texttt{@Size}) bloquent les données invalides avant tout traitement.

    \item \textbf{Gestion centralisée des erreurs}~: le \texttt{GlobalExceptionHandler} intercepte
          toutes les erreurs de validation et les exceptions métier pour retourner des messages
          lisibles et structurés au format JSON.

    \item \textbf{CORS unifié}~: la configuration CORS est centralisée dans \texttt{CorsConfig},
          évitant les conflits entre annotations et configuration globale.
\end{enumerate}

% ─────────────────────────────────────────────────────────────────────────────
\section{Conclusion}
% ─────────────────────────────────────────────────────────────────────────────

Les tests effectués confirment que LudoHeritage répond correctement à l'ensemble des scénarios
identifiés lors de l'analyse des besoins. L'application se comporte de manière robuste, même
dans des conditions dégradées (MongoDB ou Ollama hors ligne). Le chapitre suivant présente
les perspectives d'amélioration envisagées pour la plateforme.
